\documentclass[../../main.tex]{subfiles}

\begin{document}

\chapter{Analisis Algoritma Block Cipher: AES}

\begin{learningoutcome}
Setelah mempelajari bab ini, mahasiswa diharapkan mampu:
\begin{itemize}
    \item Menjelaskan latar belakang pengembangan \textit{Advanced Encryption Standard} (AES) sebagai pengganti DES (Sub-CPMK 1.2).
    \item Menganalisis spesifikasi teknis AES, termasuk varian panjang kunci (128, 192, dan 256 bit) (Sub-CPMK 1.2).
    \item Menguraikan mekanisme kerja internal AES, termasuk transformasi \textit{SubBytes}, \textit{ShiftRows}, \textit{MixColumns}, dan \textit{AddRoundKey} (Sub-CPMK 1.2).
    \item Memahami konsep \textit{state} dan representasi data dalam matriks $4 \times 4$ pada AES (Sub-CPMK 1.2).
    \item Menjelaskan proses pembangkitan kunci putaran (\textit{Key Expansion}) pada AES (Sub-CPMK 1.2).
    \item Memahami dasar matematika $\text{GF}(2^8)$ yang digunakan dalam operasi AES (Sub-CPMK 1.2).
\end{itemize}
\end{learningoutcome}

\subfile{section-01}
\subfile{section-02}
\subfile{section-03}
\subfile{section-04}
\section{Contoh Terhitung}
Nilai pada contoh berikut diverifikasi dengan perhitungan Python (S-box AES dan aritmetika $\text{GF}(2^8)$).

\begin{example}
\textbf{SubBytes dan ShiftRows.} Mulai dari state 4$\times$4 byte:
\[
\text{state}=\begin{bmatrix}00&01&02&03\\04&05&06&07\\08&09&0A&0B\\0C&0D&0E&0F\end{bmatrix}.
\]
\textit{SubBytes} mengganti setiap byte $xy$ memakai S-box pada baris $x$ kolom $y$:
\[
\begin{bmatrix}00&01&02&03\\04&05&06&07\\08&09&0A&0B\\0C&0D&0E&0F\end{bmatrix}
\longrightarrow
\begin{bmatrix}63&7C&77&7B\\F2&6B&6F&C5\\30&01&67&2B\\FE&D7&AB&76\end{bmatrix},
\]
misalnya $00\to\code{63}$ (kolom 0 baris 0) dan $53\to\code{ED}$. \textit{ShiftRows} kemudian menggeser siklik baris ke-$r$ ke kiri sejauh $r$ byte ($r=0,1,2,3$):
\[
\begin{bmatrix}63&7C&77&7B\\F2&6B&6F&C5\\30&01&67&2B\\FE&D7&AB&76\end{bmatrix}
\longrightarrow
\begin{bmatrix}63&7C&77&7B\\6B&6F&C5&F2\\67&2B&30&01\\76&FE&D7&AB\end{bmatrix}.
\]
Jadi baris 0 tetap, sedangkan baris 1, 2, dan 3 bergeser 1, 2, dan 3 byte.
\end{example}

\begin{example}
\textbf{MixColumns di $\text{GF}(2^8)$.} Satu kolom state dikalikan matriks MixColumns:
\[
\begin{bmatrix}02&03&01&01\\01&02&03&01\\01&01&02&03\\03&01&01&02\end{bmatrix}
\begin{bmatrix}26\\7B\\BD\\43\end{bmatrix}
=
\begin{bmatrix}3F\\4F\\F9\\2A\end{bmatrix}.
\]
Tinjau baris pertama: $(02\bullet26)\oplus(03\bullet7B)\oplus(01\bullet BD)\oplus(01\bullet43)$. Dalam $\text{GF}(2^8)$, $02\bullet26$ adalah geser kiri 1 bit tanpa reduksi, yaitu $\code{4C}$, dan
\[
03\bullet7B=(02\bullet7B)\oplus7B=\code{F6}\oplus\code{7B}=\code{8D}.
\]
Maka
\[
\code{4C}\oplus\code{8D}=\code{C1},\qquad \code{C1}\oplus\code{BD}=\code{7C},\qquad \code{7C}\oplus\code{43}=\code{3F},
\]
sesuai elemen pertama hasil, sehingga seluruh kolom menjadi $[\,3F\ 4F\ F9\ 2A\,]^T$.
\end{example}


\begin{obeactivity}
\textbf{Aktivitas 8.1: Simulasi Transformasi AES}
1. Ambil sebuah blok data 16-byte (misal: "Kriptografi-AES!").
2. Representasikan data tersebut ke dalam matriks \texttt{state} $4 \times 4$.
3. Lakukan simulasi tahap \textit{ShiftRows} secara manual pada matriks tersebut.
4. Amati bagaimana posisi byte berubah dan diskusikan pengaruhnya terhadap difusi data.
\end{obeactivity}

Langkah-langkah pada Aktivitas 8.1 dapat dijalankan langsung melalui program
pada Listing~\ref{lst:aes-shiftrows}. Program itu memetakan 16 byte ke matriks
\texttt{state}, menerapkan \texttt{ShiftRows} beserta kebalikannya, lalu
menunjukkan perpindahan posisi setiap byte pada matriks.

\lstinputlisting[
    language=Python,
    caption={Transformasi ShiftRows: representasi state dan difusi antar kolom},
    label={lst:aes-shiftrows}
]{\codefile{python/aes_shiftrows.py}}

Program yang sama tersedia dalam C++ (\texttt{code/cpp/aes\_shiftrows.cpp}) dan
MATLAB (\texttt{code/matlab/aes\_shiftrows.m}).

\begin{obereflection}
\textbf{Latihan 8.1}
1. Mengapa tahap \textit{MixColumns} dihilangkan pada putaran terakhir AES?
2. Jelaskan bagaimana kombinasi \textit{SubBytes}, \textit{ShiftRows}, dan \textit{MixColumns} secara kolektif mengimplementasikan prinsip \textit{confusion} dan \textit{diffusion} Shannon.
3. Berikan analisis singkat: mengapa AES-256 lebih tahan terhadap serangan \textit{quantum computing} (misal: algoritma Grover) dibandingkan AES-128?

\textbf{Latihan 8.2}
4. Jelaskan mengapa pemetaan 16 byte masukan ke matriks \textit{state} harus dilakukan
   kolom demi kolom, bukan baris demi baris. Apa akibat praktisnya bagi implementasi
   yang salah memilih konvensi?
5. Buktikan bahwa AddRoundKey bersifat involutif, sedangkan SubBytes dan MixColumns
   tidak. Untuk yang terakhir, tunjukkan dua byte berbeda yang dipertukarkan oleh
   SubBytes.
6. Hitunglah dengan tangan $\code{02} \bullet \code{57}$ dan
   $\code{04} \bullet \code{57}$ di $\text{GF}(2^8)$, lalu tentukan
   $\code{08} \bullet \code{57}$. Petunjuk: setiap langkah hanya menggeser satu bit
   ke kiri dan mereduksi dengan $\code{1B}$ bila perlu.
7. Diketahui bahwa invers perkalian dari $\code{1F}$ di $\text{GF}(2^8)$ adalah
   $\code{B2}$. Periksa dengan Persamaan~\eqref{eq:aes-sbox-rumus} bahwa
   $\text{S-box}(\code{1F}) = \code{C0}$, lalu periksa hasil itu pada
   Tabel~\ref{tab:sbox-aes}.
8. Bangkitkan $\text{RC}[11]$ seandainya AES memerlukan sebelas kunci putaran
   untuk varian berputaran 11. Petunjuk: lanjutkan dari
   $\text{RC}[10] = \code{36}$.

\textbf{Latihan 8.3}
9. Tuliskan urutan transformasi yang dijalankan oleh satu putaran dekripsi AES dan
   bandingkan dengan urutan pada putaran enkripsi. Jelaskan mengapa keduanya tidak
   sekadar kebalikan cermin.
10. Pada \textit{equivalent inverse cipher}, kunci putaran diubah menjadi
    $\tilde{K}_i = \text{InvMixColumns}(K_i)$. Tunjukkan dengan mempergunakan
    sifat linier InvMixColumns bahwa perubahan itu tidak mengubah hasil dekripsi.
11. Tunjukkan bahwa InvShiftRows dan InvSubBytes saling komutatif, tetapi
    InvMixColumns tidak komutatif dengan keduanya. Berikan satu contoh penyangkal
    untuk yang terakhir.
12. Untuk AES-256, jelaskan mengapa aturan $\text{SubWord}$ tambahan diperlukan
    hanya ketika $i \bmod 8 = 4$, dan apa yang terjadi pada jadwal kunci bila aturan
    itu dihilangkan.
13. Berapa banyak word yang harus dibangkitkan untuk AES-192, dan berapa nilai
    $\text{Rcon}$ terbesar yang dipakai? Jelaskan asal angka-angka itu.
14. AES-128 memerlukan 10 putaran, sedangkan DES 16 putaran atas blok yang justru
    lebih kecil. Jelaskan mengapa AES tetap dapat mencapai difusi penuh lebih cepat.
\end{obereflection}

\section{Rangkuman}
\begin{itemize}
    \item AES adalah algoritma Rijndael (Rijmen dan Daemen) yang dipilih NIST pada 2001: blok tetap 128 bit dengan kunci 128, 192, atau 256 bit untuk 10, 12, atau 14 putaran.
    \item Berbeda dari DES yang memakai jaringan Feistel, AES memakai struktur substitusi--permutasi (SPN) yang beroperasi pada \textit{state} matriks $4\times4$ byte.
    \item Setiap putaran utama menjalankan SubBytes, ShiftRows, MixColumns, dan AddRoundKey; putaran terakhir sama tetapi tanpa MixColumns.
    \item SubBytes memberikan \textit{confusion} melalui S-box non-linear yang dibangun dari invers di $\text{GF}(2^8)$ dan transformasi affine; ShiftRows memberikan difusi antar kolom dengan pergeseran siklik baris.
    \item MixColumns memperkuat difusi dengan mengalikan state memakai matriks tetap $\{02,03,01,01\}$ di $\text{GF}(2^8)$, sedangkan AddRoundKey hanya men-XOR state dengan kunci putaran.
    \item KeyExpansion membangkitkan kunci putaran dari kunci eksternal melalui RotWord, SubWord, dan konstanta putaran Rcon.
    \item Aritmetika AES berada di $\text{GF}(2^8)$: byte direpresentasikan sebagai polinomial derajat 7, penjumlahan adalah XOR, dan perkalian direduksi oleh polinomial $x^8+x^4+x^3+x+1$.
\end{itemize}


\begin{obeassessment}
\textbf{Kuis Bab 8}
1. Apa perbedaan utama antara struktur internal DES (Feistel) dan AES (Substitusi-Permutasi)?
2. Jelaskan peran \textit{S-box} dalam AES dan bagaimana ia dibentuk secara matematis.
3. Mengapa operasi dalam AES dilakukan dalam $\text{GF}(2^8)$ dan bukan aritmetika bilangan bulat biasa?
4. Tuliskan urutan transformasi pada satu putaran dekripsi AES dan jelaskan mengapa
   urutan itu tidak dapat diperoleh hanya dengan membalik urutan transformasi pada
   enkripsi.
5. Sebuah implementasi mandiri menghasilkan cipherteks yang berbeda dari
   implementasi lain untuk plainteks dan kunci yang sama. Sebutkan dua penyebab
   yang paling mungkin dan cara memeriksanya.
6. Diketahui state pada akhir putaran ke-9 sama dengan
   \[
   \code{EB40F21E592E38848BA113E71BC342D2},
   \]
   dan kunci putaran terakhir $K_{10}$ sama dengan
   \[
   \code{A150BE1DEFC16E55C4C018A35610A1DD}.
   \]
   Hitunglah cipherteksnya, dengan memperhatikan bahwa putaran terakhir tidak
   memuat MixColumns.
7. Bangkitkan $w[8]$ beserta $K_1$ untuk kunci AES-256 yang tersusun atas delapan
   word berikut:
   \[
   \code{603DEB10}\ \code{15CA71BE}\ \code{2B73AEF0}\ \code{857D7781}
   \]
   \[
   \code{1F352C07}\ \code{3B6108D7}\ \code{2D9810A3}\ \code{0914DFF4}.
   \]
   Jelaskan aturan mana yang Anda pakai, dan tunjukkan pada word keberapakah
   aturan khusus AES-256 mulai berlaku.
\end{obeassessment}
\begin{competencychecklist}
\begin{tabularx}{\linewidth}{|X|c|}
\hline
Indikator Kompetensi & Tercapai \\
\hline
Saya dapat menjelaskan alasan penggantian DES menjadi AES & [ ] \\
\hline
Saya dapat membedakan varian AES-128, 192, dan 256 & [ ] \\
\hline
Saya dapat menguraikan alur kerja \textit{SubBytes, ShiftRows, MixColumns} & [ ] \\
\hline
Saya dapat menjelaskan konsep matriks \textit{state} pada AES & [ ] \\
\hline
Saya dapat memahami dasar operasi aritmetika pada $\text{GF}(2^8)$ & [ ] \\
\hline
Saya dapat menjelaskan urutan transformasi pada dekripsi AES & [ ] \\
\hline
Saya dapat membedakan \textit{straightforward} dan \textit{equivalent inverse cipher} & [ ] \\
\hline
Saya dapat membangkitkan S-box AES dari konstruksi $\text{GF}(2^8)$ & [ ] \\
\hline
Saya dapat menghitung kunci putaran AES-128 dan AES-256 dengan tangan & [ ] \\
\hline
Saya dapat menunjukkan bahwa implementasi AES saya cocok dengan vektor uji FIPS 197 & [ ] \\
\hline
\end{tabularx}
\end{competencychecklist}


\end{document}
